\section{Methodology}
\label{sec:methodology}

Our insight --- that API ecosystem selection bias, not implementation error, explains metadata coverage gaps --- demands a pipeline design that can distinguish between the two. This section describes the infrastructure feasibility study (H-E1) that implements and executes this distinction.

\subsection{Study Design}

We design H-E1 as a \emph{gate hypothesis}: a pre-specified infrastructure feasibility check with binary outcome that must pass before any regression analysis proceeds. This design is motivated by the risk of running statistical analysis on severely incomplete data --- with 30\% independent variable coverage, any regression coefficients would be confounded with the domain distribution of datasets with cards, not the domain distribution of datasets in Raff's corpus.

The gate criteria are pre-specified (not chosen post-hoc):
\begin{itemize}
  \item \textbf{HF Hub coverage threshold:} $\geq$50\% of Raff's unique datasets must return valid dataset cards
  \item \textbf{OpenML temporal filter threshold:} $\geq$70\% of Raff's unique datasets must have pre-publication OpenML entries valid for Herfindahl-Hirschman Index computation
\end{itemize}

These thresholds ensure that the independent variables (documentation completeness IV1, concentration IV2) are computable for a sufficient fraction of the 255-paper corpus to support valid logistic regression without severe survivorship bias.

\subsection{Pipeline Architecture}

The H-E1 pipeline consists of five modules orchestrated by a single entry-point script (\texttt{pipeline.py}):

\textbf{RaffParser} loads Raff's CSV (255 papers, binary reproducibility labels, paper-quality features) and extracts the set of unique benchmark datasets ($N=50$) with their earliest paper publication years. The parser validates that all expected columns are present and that 255 rows are loaded, raising on violation.

\textbf{HFCoverageChecker} queries HuggingFace Hub for each of the 50 datasets using \texttt{DatasetCard.load(name)} with canonical dataset names (e.g., ``cifar10'', ``imagenet-1k'', ``coco''). For each returned card, it computes a field-presence score as the fraction of 7 Datasheets for Datasets schema fields present: \{\texttt{intended\_use}, \texttt{out\_of\_scope\_use}, \texttt{limitations}, \texttt{license}, \texttt{task\_categories}, \texttt{dataset\_info}, \texttt{provenance}\}. Rate limiting (1.0s sleep between requests) prevents API throttling. HTTP 403 errors (authentication required) and 404 errors (dataset not found) are explicitly handled and distinguished in the output.

\textbf{Rationale for canonical names:} A reproducibility auditing pipeline must be automatable. Canonical dataset names are the correct interface for a scalable pipeline; requiring name variant exploration defeats the purpose of automated auditing.

\textbf{OpenMLTemporalChecker} performs a single bulk API call (\texttt{openml.datasets.list\_datasets(output\_format=`dataframe')}) to retrieve all available OpenML datasets with metadata including upload date. It then applies a client-side temporal filter: for each of the 50 Raff datasets, it checks whether a matching OpenML entry exists with \texttt{upload\_date\_year} $<$ \texttt{paper\_publication\_year}. This filter ensures pre-publication run counts (the HHI proxy) are computable without contaminating the measurement with post-publication community use.

\textbf{MetricsAggregator} computes the two gate metrics from checker outputs and evaluates pass/fail against pre-specified thresholds. It additionally verifies \emph{pipeline activation}: all four activation indicators must be true (HF API reachable, OpenML API reachable, Raff CSV parseable, temporal filter executable) before gate pass/fail is meaningful.

\textbf{Visualizer} generates three figures: a gate metrics bar chart (Figure~\ref{fig:gate}), an OpenML pre-publication run count distribution histogram (Figure~\ref{fig:openml}), and an HF card field-presence heatmap (Figure~\ref{fig:heatmap}).

\subsection{Test Suite}

An 18-test pytest suite covers all pipeline components, run against the live pipeline (no mocking). Tests verify: correct Raff CSV parsing (255 rows, all columns), HF Hub card loading and field scoring, OpenML temporal filtering, gate evaluation, and figure generation.

\textbf{Rationale for no mocking:} Mocking the HF Hub or OpenML API validates the pipeline implementation but cannot detect platform scope limitations --- precisely the discovery we are trying to make. Live API tests are the only valid approach for an infrastructure feasibility study.

\subsection{Pre-Specified Gate Logic}

The gate uses AND logic: \emph{both} HF coverage rate $\geq$50\% AND OpenML temporal filter success rate $\geq$70\% must hold. A partial pass would indicate partial infrastructure adequacy; the downstream mechanism hypotheses (H-M1: documentation completeness $\to$ reproducibility; H-M2: concentration HHI $\to$ reproducibility; H-M3: field-level importance ordering) require both independent variables to be computable. If H-E1 gate fails, all downstream hypotheses are blocked.
